\documentclass[opre]{informs3}
\OneAndAHalfSpacedXI

\usepackage[utf8]{inputenc}
\usepackage{amsmath,amssymb,mathrsfs}
\usepackage{bm,bbm,booktabs,graphicx}
\usepackage{natbib}
\bibpunct[, ]{(}{)}{,}{a}{}{,}%
\def\bibfont{\small}%
\def\bibsep{\smallskipamount}%
\def\bibhang{24pt}%
\def\newblock{\ }%
\def\BIBand{and}%
\usepackage[hyperfootnotes=false,colorlinks=true,linkcolor=black,urlcolor=black,citecolor=blue]{hyperref}
\usepackage{cleveref}

\TheoremsNumberedThrough
\EquationsNumberedThrough
\allowdisplaybreaks

\newcommand{\E}{\mathbb{E}}
\newcommand{\Pb}{\mathbb{P}}
\newcommand{\Natu}{\mathbb{N}}
\newcommand{\Ind}[1]{\mathbbm{1}\{#1\}}
\newcommand{\State}{\mathcal{S}}
\newcommand{\Xscr}{\mathcal{X}}

\begin{document}

\RUNAUTHOR{Natalia Trigo}
\RUNTITLE{Average-Cost Dynamic Reputation Model}

\TITLE{An Average-Cost Dynamic Reputation Model with Hidden Product Choice}

\ARTICLEAUTHORS{
\AUTHOR{Natalia Trigo}
\AFF{Anderson UCLA, \EMAIL{natalia.trigo.phd@anderson.ucla.edu}}
}

\ABSTRACT{
We study an infinite-horizon reputation problem in which a seller privately chooses
between a low-cost product and a high-quality product while users learn only from
realized binary outcomes. The user treats the seller as a single Bernoulli arm and
uses Thompson sampling to decide whether to engage with the seller or with an
outside option. The seller chooses products to maximize long-run average profit
per calendar period. The resulting control problem is an average-reward Markov
decision process on the user's posterior sufficient statistics. The formulation
characterizes the seller's optimal hidden product choice through the relative
value of a success versus a failure.
}

\KEYWORDS{dynamic reputation; hidden actions; average-cost Markov decision processes; Thompson sampling}

\maketitle

\section{Model Formulation}

Consider an infinite-horizon discrete-time setting, indexed by
$t\in\Natu^+=\{1,2,\ldots\}$. The environment consists of a single user and two
sellers, denoted by $A$ and $B$. At each period $t$, the user selects a seller,
denoted by $y_t\in\{A,B\}$. The user's realized utility in period $t$ is a binary
random variable $u_t\in\{0,1\}$.

Seller B offers a transparent outside option. If the user selects Seller B, then
the realized utility is drawn from a Bernoulli distribution with known success
probability $p_0\in(0,1)$:
\begin{equation}
    \Pb(u_t=1\mid y_t=B)=p_0.
    \label{eq:seller-b-success}
\end{equation}

Seller A has access to two products, indexed by
$\Xscr=\{1,2\}$. When the user selects Seller A, the seller privately chooses a
product $x_t\in\Xscr$. Product $i\in\Xscr$ generates a successful outcome for the
user with probability $p_i$ and imposes cost $c_i$ on Seller A:
\begin{equation}
    \Pb(u_t=1\mid y_t=A,x_t=i)=p_i.
    \label{eq:seller-a-success}
\end{equation}
We assume
\begin{equation}
    0<p_1<p_2\leq 1,\qquad c_1<c_2,
    \label{eq:product-ordering}
\end{equation}
so product 2 is more likely to satisfy the user but is also more costly. Whenever
Seller A is chosen, it receives a fixed revenue $R>0$. Seller A's period profit is
\begin{equation}
    \pi_t^A=
    \begin{cases}
        R-c_{x_t} & \text{if } y_t=A,\\
        0 & \text{if } y_t=B.
    \end{cases}
    \label{eq:period-profit}
\end{equation}

\section{Information Structure and Learning Dynamics}

The central feature of the model is that Seller A's product choice is hidden from
the user. The product affects the distribution of realized utility, but the user
does not observe which product generated the realized outcome.

\begin{assumption}[User's Misspecified Belief]
The user observes the realized utility $u_t$ after choosing Seller A, but does
not observe the product $x_t\in\Xscr$ assigned by Seller A. The user therefore
treats Seller A as a single product characterized by an unknown Bernoulli success
probability $\theta\in[0,1]$.
\end{assumption}

The user learns about $\theta$ through Bayesian updating. Let $S_t$ and $F_t$
denote the cumulative number of successes and failures generated by Seller A up
to the beginning of period $t$:
\begin{align}
    S_t &= \sum_{\tau=1}^{t-1}\Ind{y_\tau=A,u_\tau=1},\\
    F_t &= \sum_{\tau=1}^{t-1}\Ind{y_\tau=A,u_\tau=0}.
\end{align}
The user starts from the uniform prior $\theta\sim\operatorname{Beta}(1,1)$.
Given state $(S_t,F_t)$, the posterior belief is
\begin{equation}
    \theta\mid S_t,F_t\sim\operatorname{Beta}(1+S_t,1+F_t).
    \label{eq:posterior}
\end{equation}

The user follows a Thompson sampling heuristic. At the beginning of period $t$,
the user samples
\begin{equation}
    \tilde{\theta}_t\sim\operatorname{Beta}(1+S_t,1+F_t).
\end{equation}
Because the outside option has known success probability $p_0$, the user chooses
Seller A if and only if the sampled index exceeds $p_0$:
\begin{equation}
    y_t=
    \begin{cases}
        A & \text{if } \tilde{\theta}_t\geq p_0,\\
        B & \text{if } \tilde{\theta}_t<p_0.
    \end{cases}
    \label{eq:user-policy}
\end{equation}

Seller A observes the user's choices and realized outcomes, and can therefore
track the sufficient statistics $(S_t,F_t)$. For a state $(S,F)$, define
Seller A's demand probability by
\begin{equation}
    \rho(S,F)
    \triangleq
    \Pb(\tilde{\theta}_t\geq p_0\mid S_t=S,F_t=F)
    =
    \int_{p_0}^{1}
    \frac{\theta^S(1-\theta)^F}{B(1+S,1+F)}\,d\theta,
    \label{eq:demand-probability}
\end{equation}
where $B(\cdot,\cdot)$ is the Beta function.

\section{Seller A's Average-Cost Optimization Problem}

Seller A faces a Markov decision problem with state
$\bm{s}_t=(S_t,F_t)\in\Natu_0\times\Natu_0$. A stationary policy is a mapping
$\pi_A:\Natu_0\times\Natu_0\to\Xscr$ that assigns a hidden product to each
posterior state. The seller's objective is to maximize long-run average profit
per calendar period:
\begin{equation}
    \eta^*
    =
    \sup_{\pi_A}
    \liminf_{T\to\infty}
    \frac{1}{T}
    \E_{\pi_A}
    \left[
        \sum_{t=1}^{T}\Ind{y_t=A}(R-c_{x_t})
    \right].
    \label{eq:average-profit}
\end{equation}
Equivalently, the same problem can be written as an average-cost problem by
minimizing the negative of Seller A's profit.

Let $\eta$ denote the optimal average profit and let $h(S,F)$ denote the relative
value, or bias, of state $(S,F)$. When the user chooses Seller B, calendar time
advances but Seller A's reputation state remains $(S,F)$. When the user chooses
Seller A and product $x$ is used, the next state is $(S+1,F)$ after a success and
$(S,F+1)$ after a failure. The average-reward optimality equation is therefore
\begin{equation}
\begin{aligned}
    \eta+h(S,F)
    =
    &[1-\rho(S,F)]h(S,F)\\
    &+\rho(S,F)
    \max_{x\in\{1,2\}}
    \left\{
        R-c_x
        +p_xh(S+1,F)
        +(1-p_x)h(S,F+1)
    \right\}.
\end{aligned}
\label{eq:average-bellman}
\end{equation}

For each product $x$, define the product-specific average-reward action value
\begin{equation}
    Q_x(S,F)
    \triangleq
    R-c_x
    +p_xh(S+1,F)
    +(1-p_x)h(S,F+1).
    \label{eq:average-action-value}
\end{equation}
The demand probability $\rho(S,F)$ determines how often Seller A has the
opportunity to act, while the product choice at a fixed state is governed by the
relative values embedded in $Q_x(S,F)$. Seller A chooses product 2 whenever
\begin{equation}
    Q_2(S,F)\geq Q_1(S,F),
\end{equation}
or equivalently whenever
\begin{equation}
    h(S+1,F)-h(S,F+1)
    \geq
    \frac{c_2-c_1}{p_2-p_1}.
    \label{eq:average-threshold}
\end{equation}
Thus, the high-quality product is optimal precisely when the relative value of a
success rather than a failure is large enough to justify the additional product
cost.

% \section{Finite-State Approximation}

% The state space in \cref{eq:average-bellman} is unbounded because each Seller A
% interaction adds one observation to the user's posterior history. The numerical
% implementation solves a finite approximation on the triangular grid
% \begin{equation}
%     \State_N=\{(S,F)\in\Natu_0^2:S+F\leq N\}.
% \end{equation}
% For states with $S+F<N$, the transition after a Seller A interaction follows the
% Bayesian sufficient-statistic update:
% \begin{equation}
%     (S,F)\mapsto(S+1,F)
%     \quad\text{after a success,}\qquad
%     (S,F)\mapsto(S,F+1)
%     \quad\text{after a failure.}
% \end{equation}

% At the boundary $S+F=N$, the approximation uses a rolling observation window. A
% new outcome is added, and one of the existing $N$ observations is dropped
% uniformly from the empirical history. After a success,
% \begin{equation}
%     (S,F)\mapsto
%     \begin{cases}
%         (S+1,F-1) & \text{with probability } F/N,\\
%         (S,F) & \text{with probability } S/N,
%     \end{cases}
%     \label{eq:boundary-success}
% \end{equation}
% and after a failure,
% \begin{equation}
%     (S,F)\mapsto
%     \begin{cases}
%         (S-1,F+1) & \text{with probability } S/N,\\
%         (S,F) & \text{with probability } F/N.
%     \end{cases}
%     \label{eq:boundary-failure}
% \end{equation}
% Transitions with zero probability are ignored at the corners of the boundary.
% This construction avoids absorbing boundary classes while updating the empirical
% success share in the direction of the newly observed outcome.

% For a relative value function $h$ on $\State_N$, let $\bar h^+(S,F)$ and
% $\bar h^-(S,F)$ denote the expected next relative value after a success and after
% a failure, respectively, under the finite-state transition rule. For interior
% states, $\bar h^+(S,F)=h(S+1,F)$ and $\bar h^-(S,F)=h(S,F+1)$; at boundary states,
% these quantities are computed using \cref{eq:boundary-success,eq:boundary-failure}.
% The finite-grid counterpart of \cref{eq:average-threshold} is
% \begin{equation}
%     \bar h^+(S,F)-\bar h^-(S,F)
%     \geq
%     \frac{c_2-c_1}{p_2-p_1}.
%     \label{eq:finite-average-threshold}
% \end{equation}

% \section{Average-Reward Policy Iteration}

% The finite-state problem is solved by average-reward policy iteration. Fix a
% stationary policy $\pi$ on $\State_N$. Let $P_\pi^A$ be the transition matrix over
% posterior states conditional on the user choosing Seller A, and let
% $r_\pi(S,F)=R-c_{\pi(S,F)}$ be the profit conditional on Seller A being chosen.
% The policy evaluation equations are
% \begin{equation}
%     \eta_\pi+\rho(S,F)h_\pi(S,F)
%     =
%     \rho(S,F)
%     \left[
%         r_\pi(S,F)
%         +
%         \sum_{(S',F')\in\State_N}
%         P_\pi^A((S,F),(S',F'))h_\pi(S',F')
%     \right],
%     \label{eq:policy-evaluation}
% \end{equation}
% with normalization $h_\pi(0,0)=0$. This form is obtained from
% \cref{eq:average-bellman} by collecting the self-loop term that occurs when the
% user chooses Seller B.

% Given the evaluated bias $h_\pi$, policy improvement chooses
% \begin{equation}
%     \pi^+(S,F)\in
%     \arg\max_{x\in\{1,2\}}
%     \left\{
%         R-c_x
%         +
%         \sum_{(S',F')\in\State_N}
%         P_x^A((S,F),(S',F'))h_\pi(S',F')
%     \right\},
%     \label{eq:policy-improvement}
% \end{equation}
% where $P_x^A$ is the conditional transition matrix induced by product $x$. The
% iteration terminates when the improved policy no longer changes. At that point,
% the policy satisfies the finite-state average-reward optimality equation.

% \section{Numerical Study}

% The numerical study evaluates how the optimal hidden product policy depends on
% the outside option quality $p_0$. For each value of $p_0$, the solver computes
% the average-profit gain $\eta$, the relative value function $h$, the optimal
% product choice over posterior states, and the finite-grid bias gap
% $\bar h^+(S,F)-\bar h^-(S,F)$ that determines the product-2 region in
% \cref{eq:finite-average-threshold}. Simulations then generate sample histories
% under the computed stationary policy, recording Seller A's market share, product
% mix, realized average profit, posterior beliefs, and demand probabilities over
% time.

\end{document}
